\begin{proof}

\begin{enumerate}[label={$\langle1\rangle$\arabic*.},leftmargin=*]
\item Write $H$ for the horizontal section, $g=g(H)\ge1$, and
$C_H$ for its image in the reduced conductor $C$. The preceding
ruled reduction gives two rational normalized components of $C$
and one point-identification cycle. The nonconstant map
\[
                         H\longrightarrow C_H^\nu=\mathbb P^1
\]
has degree at least two: a degree-one finite map between smooth
curves is an isomorphism, whereas their genera differ. Each
generic point of $H$ has conductor exponent two, by the actual
pole divisor of P4R.6. Thus the generic transverse plane curve
along $C_H$ has at least two normalization branches of exponent two.

\item For its branches $B_i$ the exact conductor formula is
\[
             c_i=2\delta(B_i)+\sum_{j\ne i}I(B_i,B_j).
\]
A singular branch with $c_i=2$ permits no other branch,
so the two specified branches are smooth. With two total
branches their intersection multiplicity is two, giving $A_3$.
With three total branches, the first two exponents force all
three pairwise intersection multiplicities to equal one;
the third branch is then smooth as well, giving $D_4$.
Four branches would give exponent at least three on every
smooth branch. These possibilities are exhaustive. In type
$A_3$, $H$ has degree exactly two and supplies both generic
normalization preimages. In type $D_4$ its degree is two or
three. We exclude the former transverse type next.

\item Assume that the transverse germ along $C_H$ is $A_3$.
Let $C_{H,\rm sch}$ be the primary component of the actual
image conductor scheme supported on $C_H$. The primary
restriction construction of P4R.13 applies before its
elliptic-genus calculation, without change. More explicitly,
P4.33 gives rational singularities of the normalization when
$H$ is not contracted, so finite duality and the canonical
trace give the injections
\[
 \mathcal O_{C_{\rm sch}}\longrightarrow
 \nu_*(\mathcal O_N/I_N)\longrightarrow\rho_*\mathcal O_D.
\]
The exact blowup relation $D'=\beta^*D''-E$ identifies
global functions on $D$ with those on $2H_0+V_0$ of the
minimal ruled model. The kernel of restriction to $2H_0$
has successive vertical-fibre layers
$\mathcal O_{\mathbb P^1}(-2)$, with no global sections.
Hence a global function is determined at the generic doubled
horizontal section. It follows that restriction from
$C_{\rm sch}$ to $C_{H,\rm sch}$ is injective on $H^0$
and surjective on $H^1$: the first assertion follows from
those injections and generic determination, and the second
from the dimension-one kernel of the restriction sheaf map.
Therefore
\[
                 g-1=\chi(\mathcal O_{C_{\rm sch}})
                    \le\chi(\mathcal O_{C_{H,\rm sch}}).
                                                            \tag{P4R.18.1}
\]
Every arrow here is the primary restriction of the same
conductor algebra; no other conductor component is deleted
from the left side.

\item We evaluate the right side on $C_H^\nu=\mathbb P^1$.
Use the actual differentiated conormal, the saturated Hessian
line and the rank-one nilradical of P4R.10. The kernel $K$
of the normalized derivative has degree $2+\deg R$, since
$c_1(Y)\cdot C_H=0$. The Hessian line has degree
$\tfrac12\sum d_p$. Riemann--Roch and the normalization
sequence give exactly
\[
 \chi(\mathcal O_{C_{H,\rm sch}})
 =4-\sum_p\left[\delta(C_H,p)
      +\sum_{i\mapsto p}(d_i/2+b_i-r_i)+q_p\right].
                                                            \tag{P4R.18.2}
\]
All terms are those of this primary curve and its primary
conductor. They remain defined at its intersections with
the other component, because the pure conductor quotient
extends its generic containment of the symbolic square.

\item Every singular point of $C_H$ contributes at least three
to the sum in (P4R.18.2). At a planar point the ideal is
$(u,h)$, its symbolic square equals its ordinary square,
and the actual equation belongs to that square. The
rank-one symmetric-square estimate of P4.42 and the exact
plane-curve different calculation give a contribution at
least $3\delta(C_H,p)+q_p\ge3$.
At a nonplanar point every defining equation has order at
least two. On a branch of multiplicity $m_i$ its projected
conormal order is at least $m_i$. The point is singular
on the entire critical curve, so P4R.15 makes the ambient
Hessian zero and gives $d_i\ge m_i$. Since $r_i=m_i-1$,
the contribution is at least
\[
                      \delta(C_H,p)+\sum_i(1+m_i/2),
\]
which is at least three at a singular curve point. The
same arguments give nonnegative contributions everywhere;
at a smooth primary point its conormal is the full kernel
and its only possible contribution is the nonnegative
Hessian zero order divided by two.
Combining (P4R.18.1)--(P4R.18.2) gives total primary
defect at most $5-g\le4$. Thus $C_H$ has at most one
singular point, and two different singular points are
impossible even in the presence of arbitrary transverse
types on the other global component.

\item We count the possible branch values of the degree-two map
$H\to\mathbb P^1$ with the identifications retained.
The normalization graph formula is
\[
                         \sum_{p\in C}(r_p-1)=2.
\]
Connectedness leaves either two two-branch points or one
three-branch point. In the two-point case, two intersections
of the global components give two points on $C_H^\nu$.
There can be at most one additional unibranch singularity
of $C_H$, so these supply at most three possible branch
values. The other connected possibility is a self-identification
on one component and one intersection of the two components.
A self-identification on $C_H$ supplies two normalization
points and already costs at least three in the primary
defect. It permits no additional singular point, so with
the intersection there are again at most three possible
branch values. A self-identification on the other component
supplies no additional point on $C_H^\nu$.

\item At a three-branch point, the two global components must
supply two and one branches. Thus $C_H^\nu$ has at most
two points over it. Two points already make $C_H$ singular
there and exclude any further singular point by the defect
bound. One point permits at most one further unibranch
singularity. This case also has at most three possible
branch values.

\item There is no ramification elsewhere. At a point where the
entire critical curve is smooth, P4R.14 gives precisely
the product tacnode or $X(Y^2+XZ)$. Both have two distinct
normalization points over that point of the conductor.
Their two generic preimages lie on $H$, so properness
places both special preimages on $H$ as well. The
degree-two map is therefore unramified there. We have
accounted for every possible ramification value.
Riemann--Hurwitz for a degree-two map from a genus-$g$
curve requires $2g+2\ge4$ distinct branch values; each
ramified fibre consists of its single point of index two.
This contradicts the count of at most three. Thus the
generic germ along $C_H$ must be $D_4$, with degree two
or three as proved at the start.

\end{enumerate}
\end{proof}
